% Options for packages loaded elsewhere
\PassOptionsToPackage{unicode}{hyperref}
\PassOptionsToPackage{hyphens}{url}
\PassOptionsToPackage{dvipsnames,svgnames,x11names}{xcolor}
%
\documentclass[
  11pt,
  a4paper,
]{article}
\usepackage{amsmath,amssymb}
\usepackage{iftex}
\ifPDFTeX
  \usepackage[T1]{fontenc}
  \usepackage[utf8]{inputenc}
  \usepackage{textcomp} % provide euro and other symbols
\else % if luatex or xetex
  \usepackage{unicode-math} % this also loads fontspec
  \defaultfontfeatures{Scale=MatchLowercase}
  \defaultfontfeatures[\rmfamily]{Ligatures=TeX,Scale=1}
\fi
\usepackage{lmodern}
\ifPDFTeX\else
  % xetex/luatex font selection
\fi
% Use upquote if available, for straight quotes in verbatim environments
\IfFileExists{upquote.sty}{\usepackage{upquote}}{}
\IfFileExists{microtype.sty}{% use microtype if available
  \usepackage[]{microtype}
  \UseMicrotypeSet[protrusion]{basicmath} % disable protrusion for tt fonts
}{}
\makeatletter
\@ifundefined{KOMAClassName}{% if non-KOMA class
  \IfFileExists{parskip.sty}{%
    \usepackage{parskip}
  }{% else
    \setlength{\parindent}{0pt}
    \setlength{\parskip}{6pt plus 2pt minus 1pt}}
}{% if KOMA class
  \KOMAoptions{parskip=half}}
\makeatother
\usepackage{xcolor}
\usepackage[margin=24mm]{geometry}
\usepackage{longtable,booktabs,array}
\usepackage{calc} % for calculating minipage widths
% Correct order of tables after \paragraph or \subparagraph
\usepackage{etoolbox}
\makeatletter
\patchcmd\longtable{\par}{\if@noskipsec\mbox{}\fi\par}{}{}
\makeatother
% Allow footnotes in longtable head/foot
\IfFileExists{footnotehyper.sty}{\usepackage{footnotehyper}}{\usepackage{footnote}}
\makesavenoteenv{longtable}
\setlength{\emergencystretch}{3em} % prevent overfull lines
\providecommand{\tightlist}{%
  \setlength{\itemsep}{0pt}\setlength{\parskip}{0pt}}
\setcounter{secnumdepth}{5}
\ifLuaTeX
\usepackage[bidi=basic]{babel}
\else
\usepackage[bidi=default]{babel}
\fi
\babelprovide[main,import]{brazilian}
% get rid of language-specific shorthands (see #6817):
\let\LanguageShortHands\languageshorthands
\def\languageshorthands#1{}
\usepackage[T1]{fontenc}
\usepackage{lmodern}
\usepackage{microtype}
\usepackage{xcolor}
\usepackage{fancyhdr}
\usepackage{amsmath,amssymb}
\usepackage{xurl}
\definecolor{qddmink}{HTML}{17324D}
\pagestyle{fancy}
\fancyhf{}
\fancyhead[L]{\small\textcolor{qddmink}{Quad-Dimensional DataMatrix}}
\fancyhead[R]{\small RfD 0.2.0}
\fancyfoot[C]{\thepage}
\setlength{\headheight}{15pt}
\setlength{\emergencystretch}{3em}

\ifLuaTeX
  \usepackage{selnolig}  % disable illegal ligatures
\fi
\IfFileExists{bookmark.sty}{\usepackage{bookmark}}{\usepackage{hyperref}}
\IfFileExists{xurl.sty}{\usepackage{xurl}}{} % add URL line breaks if available
\urlstyle{same}
\hypersetup{
  pdftitle={Quad-Dimensional DataMatrix},
  pdfauthor={João Paulo De Luca Jr; Jorge Luiz Espanhol},
  pdflang={pt-BR},
  pdfsubject={Request for Discussion: estados GIF indexados e valores
hexadecimais},
  colorlinks=true,
  linkcolor={qddmink},
  filecolor={Maroon},
  citecolor={Blue},
  urlcolor={qddmink},
  pdfcreator={LaTeX via pandoc}}

\title{Quad-Dimensional DataMatrix}
\usepackage{etoolbox}
\makeatletter
\providecommand{\subtitle}[1]{% add subtitle to \maketitle
  \apptocmd{\@title}{\par {\large #1 \par}}{}{}
}
\makeatother
\subtitle{Request for Discussion: Estados GIF Hexadecimais e Sequências
Mantidas no App}
\author{João Paulo De Luca Jr \and Jorge Luiz Espanhol}
\date{4 de novembro de 2026 · RfD 0.2.0}

\begin{document}
\maketitle
\begin{abstract}
Quad-Dimensional DataMatrix é proposto como estrutura de discussão para
dados hexadecimais representados por estados GIF indexados. O perfil de
trabalho usa um canvas de 256 por 256 pixels, 64 por 64 módulos
efetivos, até 64 entradas hexadecimais na paleta e até 64 frames. Quatro
coordenadas lógicas identificam posição horizontal, posição vertical,
estado temporal e dimension. Seu catálogo cartesiano máximo contém
16.777.216 tuplas candidatas. Uma sequência densa seleciona uma entre 64
alternativas em cada módulo e, assim, transporta até 192 KiB de payload
bruto em um ciclo de 64 frames, nominalmente oito segundos a 8 fps. Uma
sequência mantida no app oferece uma segunda interação: avançar na
cadência escolhida, pausar por seleção do usuário e copiar o código
hexadecimal do estado para um serviço receptor. O paper desenvolve o
modelo de contagem, separa os dois papéis do tempo e examina mapeamentos
um a um entre Unicode e hexadecimal como candidatos a um futuro
dicionário ou chave. Um estudo de 16 frames e um visualizador local
bilíngue acompanham a discussão. A seção Escopo e limitações concentra
as qualificações. O documento convida ao refinamento antes da proposta
de uma especificação.
\end{abstract}

\hypertarget{propuxf3sito-e-status-da-discussuxe3o}{%
\section{Propósito e status da
discussão}\label{propuxf3sito-e-status-da-discussuxe3o}}

\textbf{Quad-Dimensional DataMatrix} é um Request for Discussion (RfD)
de \textbf{João Paulo De Luca Jr e Jorge Luiz Espanhol}. Reúne um
catálogo de quatro coordenadas, transporte GIF indexado, valores
hexadecimais e seleção de estados mantidos no app.

O objetivo imediato é tornar os objetos propostos contáveis e testáveis.
O leitor deve conseguir identificar o conteúdo de um estado, como o app
o seleciona, o que o usuário copia e qual interpretação de um valor
hexadecimal está sendo proposta. Uma etapa posterior poderá desenvolver
uma especificação a partir da discussão.

O vocabulário de trabalho é:

\begin{longtable}[]{@{}ll@{}}
\toprule\noalign{}
Termo & Significado neste RfD \\
\midrule\noalign{}
\endhead
\bottomrule\noalign{}
\endlastfoot
Pixel armazenado & Um índice na imagem armazenada \\
Módulo efetivo & Bloco de 4 × 4 com índice repetido \\
Dimension & Identificador ativo da paleta hexadecimal \\
Estado & Matriz completa em um índice temporal \\
Sequência & Conjunto ordenado de até 64 estados \\
Tupla do catálogo & Tupla candidata \((x,y,t,d)\) \\
Código & Valor associado a um estado para consumo \\
Chave & Entrada de dicionário ou segredo proposto à discussão \\
\end{longtable}

Inglês é o idioma padrão do repositório e do manuscrito. Uma edição
completa em português é mantida em paralelo. As qualificações estão
concentradas em \protect\hyperlink{escopo-e-limitauxe7uxf5es}{Escopo e
limitações}; o argumento a seguir desenvolve diretamente o modelo de
trabalho.

\hypertarget{perfil-gif-de-trabalho-e-valores-hexadecimais}{%
\section{Perfil GIF de trabalho e valores
hexadecimais}\label{perfil-gif-de-trabalho-e-valores-hexadecimais}}

\hypertarget{geometria-e-entradas-ativas}{%
\subsection{Geometria e entradas
ativas}\label{geometria-e-entradas-ativas}}

O perfil atual é descrito por

\[
W_p=H_p=256,\qquad s=4,\qquad W=H=64,\qquad K\le64,\qquad F\le64.
\]

O frame armazenado possui \(W_pH_p=65.536\) pixels. Sua quantidade de
módulos independentes é

\[
WH=\frac{W_pH_p}{s^2}=4.096.
\]

O fator \(s^2=16\) repete cada índice selecionado em um módulo quadrado.
Determina a geometria de renderização, enquanto \(WH\) determina a
quantidade de escolhas independentes de módulo.

Cada entrada ativa da paleta é um valor hexadecimal de seis dígitos,
como \(\mathtt{\#12ABEF}\). Para o estado \(t\), o mapeamento é

\[
P_t:\{0,\ldots,K-1\}\longrightarrow
\{\mathtt{\#000000},\ldots,\mathtt{\#FFFFFF}\}.
\]

A relação pública é

\[
\mathrm{dimension}\;\longleftrightarrow\;
\mathrm{\text{índice da paleta}}\;\longleftrightarrow\;
\mathrm{\text{valor hexadecimal}}.
\]

O adaptador atual trata a paleta ordenada como dicionário de código. Sua
ordem pode mudar entre estados, preservando a representação hexadecimal
de cada entrada.
\href{https://www.w3.org/Graphics/GIF/spec-gif89a.txt}{GIF89a} fornece o
contêiner de tabelas indexadas e blocos de imagem.

\hypertarget{representauxe7uxe3o-e-custo-da-tabela}{%
\subsection{Representação e custo da
tabela}\label{representauxe7uxe3o-e-custo-da-tabela}}

Um valor hexadecimal de seis dígitos possui domínio inteiro de 24 bits:

\[
16^6=2^{24}=16.777.216.
\]

O alfabeto ativo tem \(K\) entradas; escolher uma entrada custa
\(\log_2K\) bits no modelo de contagem sem restrições. Com \(K=64\),
essa escolha transporta seis bits.

Uma tabela GIF com \(K\) entradas ocupa

\[
B_{\mathrm{table}}=3K\ \text{bytes}.
\]

Com 64 entradas, a tabela ocupa 192 bytes. Se cada estado fornecer sua
própria tabela, a alocação será de \(3KF\) bytes. Descrições textuais em
hexadecimal e o contêiner codificado possuem seus próprios orçamentos de
armazenamento. O modelo mantém essas grandezas observáveis
separadamente.

\hypertarget{o-catuxe1logo-de-quatro-coordenadas}{%
\section{O catálogo de quatro
coordenadas}\label{o-catuxe1logo-de-quatro-coordenadas}}

\hypertarget{catuxe1logo-cartesiano-e-estado-selecionado}{%
\subsection{Catálogo cartesiano e estado
selecionado}\label{catuxe1logo-cartesiano-e-estado-selecionado}}

Sejam

\[
X=\{0,\ldots,W-1\},\quad
Y=\{0,\ldots,H-1\},\quad
T=\{0,\ldots,F-1\},\quad
\Delta=\{0,\ldots,K-1\}.
\]

O catálogo proposto de quatro coordenadas é

\[
\Omega=X\times Y\times T\times\Delta,\qquad
|\Omega|=WHFK.
\]

Para o perfil máximo,

\[
|\Omega|=64^4=16.777.216.
\]

Uma sequência densa escolhe uma dimension em cada posição e estado
temporal:

\[
D:T\times Y\times X\longrightarrow\Delta.
\]

De maneira equivalente, define-se um tensor indicador de quatro índices

\[
A[t,y,x,d]=
\begin{cases}
1,&D[t,y,x]=d,\\
0,&\text{caso contrário}.
\end{cases}
\]

A condição de seleção é

\[
\sum_{d=0}^{K-1}A[t,y,x,d]=1.
\]

O catálogo enumera todas as tuplas candidatas; a sequência densa
seleciona uma dimension para cada uma de suas \(WHF\) posições. Isso dá
significado preciso às quatro coordenadas lógicas do projeto.

\hypertarget{um-endereuxe7o-hexadecimal-do-catuxe1logo}{%
\subsection{Um endereço hexadecimal do
catálogo}\label{um-endereuxe7o-hexadecimal-do-catuxe1logo}}

Com os quatro intervalos de coordenadas limitados a 64, um endereço do
catálogo em ordem de linhas pode ser escrito como

\[
a=(((64t+y)64+x)64+d).
\]

Seu intervalo é

\[
0\le a\le64^4-1=\mathtt{0xFFFFFF}.
\]

A inversão é

\[
d=a\bmod64,\qquad
x=\lfloor a/64\rfloor\bmod64,
\]

\[
y=\lfloor a/64^2\rfloor\bmod64,\qquad
t=\lfloor a/64^3\rfloor.
\]

Um endereço usa 24 bits, com seis bits para cada coordenada. A
implementação expõe o endereço como inteiro hexadecimal de seis dígitos.
Esse endereço é útil para discutir dicionários, eventos esparsos e
referências locais a estados.

\hypertarget{capacidade-e-taxas-nominais}{%
\section{Capacidade e taxas
nominais}\label{capacidade-e-taxas-nominais}}

\hypertarget{das-escolhas-aos-bits}{%
\subsection{Das escolhas aos bits}\label{das-escolhas-aos-bits}}

Para \(N=WHF\) módulos selecionáveis independentemente, a quantidade de
sequências densas é

\[
|\mathcal D|=K^N.
\]

Uma mensagem binária de comprimento \(L\) exige \(2^L\) codificações
distinguíveis. Portanto,

\[
2^L\le K^N
\quad\Longrightarrow\quad
L_{\max}=\lfloor WHF\log_2K\rfloor.
\]

No perfil máximo,

\[
N=64^3=262.144,\qquad
\log_2K=6,
\]

\[
L_{\max}=1.572.864\ \text{bits},
\qquad
B_{\max}=\left\lfloor\frac{L_{\max}}8\right\rfloor
=196.608\ \text{bytes}=192\ \text{KiB}.
\]

O payload bruto por estado é

\[
B_{\mathrm{state}}=\frac{64^2\cdot6}{8}
=3.072\ \text{bytes}=3\ \text{KiB}.
\]

A quantidade de sequências completas possíveis é exponencial em \(N\). O
payload binário é seu logaritmo e cresce proporcionalmente a \(F\).

\hypertarget{taxa-do-catuxe1logo-e-taxa-de-payload}{%
\subsection{Taxa do catálogo e taxa de
payload}\label{taxa-do-catuxe1logo-e-taxa-de-payload}}

Seja \(f\) a quantidade nominal de estados exibidos por segundo. Um
ciclo completo tem duração

\[
\tau=F/f.
\]

Com \(F=64\) e \(f=8\), \(\tau=8\) segundos. O orçamento de percurso do
catálogo é

\[
R_{\Omega}=\frac{WHFK}{\tau}
=WHKf=2.097.152
\ \text{tuplas candidatas/s}.
\]

Essa é a interpretação pretendida se a abreviação \textbf{cps} for usada
para o catálogo neste RfD. Duas outras taxas descrevem os dados
selecionados:

\[
R_{\mathrm{modules}}=WHf=32.768
\ \text{módulos/s},
\]

\[
R_{\mathrm{bits}}=WHf\log_2K
=196.608\ \text{bits/s}
=24\ \text{KiB/s}.
\]

\begin{longtable}[]{@{}lr@{}}
\toprule\noalign{}
Medida & Perfil máximo \\
\midrule\noalign{}
\endhead
\bottomrule\noalign{}
\endlastfoot
Candidatas de quatro coordenadas & 16.777.216 por ciclo \\
Orçamento de tuplas candidatas & 2.097.152 por segundo a 8 fps \\
Módulos selecionados & 262.144 por ciclo \\
Payload bruto & 192 KiB por ciclo \\
Taxa de payload bruto & 24 KiB/s a 8 fps \\
\end{longtable}

Para um enquadramento futuro com \(h\) bits reservados por estado e taxa
de código \(r\), um primeiro orçamento de payload é

\[
B_{\mathrm{net}}=
\left\lfloor\frac{rF(WH\log_2K-h)}8\right\rfloor,
\qquad 0<r\le1.
\]

A expressão pressupõe reservas fixas por estado, seguidas de um código
agregado com taxa \(r\). Outra posição das reservas resulta em outro
orçamento. O arranjo eventual é um item de discussão.

\hypertarget{tempo-no-contuxeainer-e-tempo-no-app}{%
\section{Tempo no contêiner e tempo no
app}\label{tempo-no-contuxeainer-e-tempo-no-app}}

\hypertarget{caduxeancia-de-transporte}{%
\subsection{Cadência de transporte}\label{caduxeancia-de-transporte}}

Os delays do GIF são armazenados em centésimos de segundo. Um período
nominal de 125 ms por estado pode ser representado por delays alternados
de 120 e 130 ms. Para uma quantidade par de frames,

\[
\sum_{i=0}^{F-1}\delta_i=125F\ \text{ms}.
\]

O protótipo quantiza os limites cumulativos:

\[
\delta_i=
10\left(
\left\lfloor\frac{(i+1)p}{10}\right\rfloor
-\left\lfloor\frac{ip}{10}\right\rfloor
\right)\ \text{ms},
\]

onde \(p\) é o período solicitado em milissegundos. Para \(p=125\) e
\(F=64\), os delays armazenados somam 8.000 ms.

\hypertarget{sequuxeancia-mantida-no-app-e-seleuxe7uxe3o}{%
\subsection{Sequência mantida no app e
seleção}\label{sequuxeancia-mantida-no-app-e-seleuxe7uxe3o}}

A interação proposta coloca a sequência no app. O usuário ou um serviço
receptor seleciona ou consome um estado. O app pode avançar em um
intervalo escolhido \(\Delta\), inclusive um estado por minuto.

Um seletor possível de estado é

\[
i(\theta)=
\left\lfloor\frac{\theta-\theta_0}{\Delta}\right\rfloor\bmod F.
\]

Se o usuário selecionar no instante \(\theta_c\), o índice congelado
será

\[
i_c=i(\theta_c).
\]

Um mapeamento escolhido de estado para código, \(g\), fornece o valor
para a área de transferência:

\[
c=g(D_{i_c}).
\]

O visualizador local de estudo usa a string hexadecimal canônica do
payload do estado como \(g\). Inicia com \(\Delta=60\) segundos. Um
ciclo de 64 estados dura, então, 64 minutos. O exemplo de transporte a 8
fps e a interação de um frame por minuto são dois usos dos mesmos
estados ordenados.

\begin{longtable}[]{@{}
  >{\raggedright\arraybackslash}p{(\columnwidth - 4\tabcolsep) * \real{0.3333}}
  >{\raggedright\arraybackslash}p{(\columnwidth - 4\tabcolsep) * \real{0.3333}}
  >{\raggedright\arraybackslash}p{(\columnwidth - 4\tabcolsep) * \real{0.3333}}@{}}
\toprule\noalign{}
\begin{minipage}[b]{\linewidth}\raggedright
Evento
\end{minipage} & \begin{minipage}[b]{\linewidth}\raggedright
Ação do app
\end{minipage} & \begin{minipage}[b]{\linewidth}\raggedright
Lado receptor
\end{minipage} \\
\midrule\noalign{}
\endhead
\bottomrule\noalign{}
\endlastfoot
Carregar & Mantém a sequência e seu mapeamento localmente & Aguarda um
código \\
Avançar & Exibe o próximo estado & Observa ou solicita seleção \\
Clicar & Interrompe o avanço e congela o estado exibido & Usuário
seleciona o estado atual \\
Copiar & Copia o código hexadecimal associado & Usuário, serviço ou site
o consome \\
Continuar & Retoma quando solicitado & Outro estado pode ser
selecionado \\
\end{longtable}

O exemplo renderiza cada estado selecionado como um GIF independente. O
clique cancela o temporizador do app antes de escrever na área de
transferência. Assim, a imagem selecionada e o código copiado se referem
ao mesmo estado.

\hypertarget{relauxe7uxe3o-com-totp}{%
\subsection{Relação com TOTP}\label{relauxe7uxe3o-com-totp}}

A comparação com
\href{https://www.rfc-editor.org/rfc/rfc6238.html}{TOTP} diz respeito à
posse do estado evolutivo e ao papel do verificador receptor. No TOTP,
um app mantém um segredo e calcula um código usando um contador
temporal. Nesta proposta, o app mantém ou obtém uma sequência e expõe o
código de um estado selecionado.

O serviço receptor é o consumidor ou verificador nos dois padrões de
interação. O RfD convida à discussão sobre a escolha de \(g\),
provisionamento, renovação dos estados e regras de aceitação para esse
modelo de sequência. A
\href{https://www.w3.org/TR/clipboard-apis/}{Clipboard API} fornece a
operação de cópia acionada pelo usuário utilizada no exemplo local.

\hypertarget{unicode-e-mapeamentos-hexadecimais-um-a-um}{%
\section{Unicode e mapeamentos hexadecimais um a
um}\label{unicode-e-mapeamentos-hexadecimais-um-a-um}}

\hypertarget{o-domuxednio-escalar}{%
\subsection{O domínio escalar}\label{o-domuxednio-escalar}}

O
\href{https://www.unicode.org/versions/Unicode18.0.0/core-spec/chapter-3/}{domínio
escalar Unicode} é

\[
\mathcal U=
\{0,\ldots,\mathtt{0x10FFFF}\}
\setminus
\{\mathtt{0xD800},\ldots,\mathtt{0xDFFF}\}.
\]

Sua cardinalidade é

\[
|\mathcal U|=
\mathtt{0x110000}-\mathtt{0x800}
=1.112.064.
\]

Os ``aproximadamente um milhão de códigos'' da proposta se relacionam
naturalmente com esse domínio.
\href{https://www.rfc-editor.org/rfc/rfc3629.html}{UTF-8} é uma
codificação de comprimento variável desses escalares, usando de um a
quatro octetos.

Uma representação hexadecimal direta é

\[
h(u)=\operatorname{hex}_6(u).
\]

Ela é injetiva: \(h(u_1)=h(u_2)\) implica \(u_1=u_2\). A representação
cabe no domínio de valores de 24 bits já usado pelo catálogo e pelas
descrições da paleta.

\hypertarget{truxeas-interpretauxe7uxf5es-candidatas}{%
\subsection{Três interpretações
candidatas}\label{truxeas-interpretauxe7uxf5es-candidatas}}

\textbf{Dicionário direto de paleta.} Um escalar pode ser associado a
uma entrada hexadecimal de seis dígitos. Um frame pode utilizar até
\(K\) entradas escalares distintas nessa interpretação. Um dicionário
adaptativo poderia alterar as entradas ativas entre estados. Usos
repetidos de uma entrada continuam como escolhas repetidas comuns de
módulo.

\textbf{Representação escalar em quatro dimensions.} Um alfabeto fixo de
64 identificadores de dimension pode representar qualquer escalar usando
quatro dígitos na base 64:

\[
u=64^3d_0+64^2d_1+64d_2+d_3.
\]

Os dígitos são

\[
d_j=
\left\lfloor\frac{u}{64^{3-j}}\right\rfloor\bmod64,
\qquad j=0,1,2,3.
\]

Como

\[
64^3<|\mathcal U|<64^4,
\]

quatro dígitos de comprimento fixo na base 64 são suficientes, enquanto
três não conseguem identificar todos os escalares. O protótipo oferece
esse mapeamento reversível.

\textbf{Dicionário de endereços do catálogo.} O inteiro do escalar
também pode ser inserido no domínio de endereços de 24 bits do catálogo.
As equações inversas de endereço produzem uma tupla de coordenadas. Isso
permite uma entrada de dicionário um a um ou um evento endereçado. Uma
sequência de eventos endereçados também precisa de uma representação da
ordem e da multiplicidade; essas passam a ser escolhas explícitas de
projeto.

\begin{longtable}[]{@{}
  >{\raggedright\arraybackslash}p{(\columnwidth - 6\tabcolsep) * \real{0.2500}}
  >{\raggedright\arraybackslash}p{(\columnwidth - 6\tabcolsep) * \real{0.2500}}
  >{\raggedright\arraybackslash}p{(\columnwidth - 6\tabcolsep) * \real{0.2500}}
  >{\raggedright\arraybackslash}p{(\columnwidth - 6\tabcolsep) * \real{0.2500}}@{}}
\toprule\noalign{}
\begin{minipage}[b]{\linewidth}\raggedright
Exemplo de escalar
\end{minipage} & \begin{minipage}[b]{\linewidth}\raggedright
Escalar hexadecimal
\end{minipage} & \begin{minipage}[b]{\linewidth}\raggedright
Octetos UTF-8
\end{minipage} & \begin{minipage}[b]{\linewidth}\raggedright
Quatro identificadores de dimension
\end{minipage} \\
\midrule\noalign{}
\endhead
\bottomrule\noalign{}
\endlastfoot
A latino maiúsculo & \texttt{0x000041} & \texttt{41} &
\texttt{00\ 00\ 01\ 01} \\
a latino minúsculo com til & \texttt{0x0000E3} & \texttt{C3\ A3} &
\texttt{00\ 00\ 03\ 23} \\
Sinal de euro & \texttt{0x0020AC} & \texttt{E2\ 82\ AC} &
\texttt{00\ 02\ 02\ 2C} \\
Rosto sorridente & \texttt{0x01F600} & \texttt{F0\ 9F\ 98\ 80} &
\texttt{00\ 1F\ 18\ 00} \\
\end{longtable}

\hypertarget{oruxe7amento-de-texto-e-eventual-chave}{%
\subsection{Orçamento de texto e eventual
chave}\label{oruxe7amento-de-texto-e-eventual-chave}}

A representação em quatro dimensions usa 24 bits por escalar:

\[
B_{\mathrm{scalar}}=3\ \text{bytes},
\qquad
C_{\mathrm{state}}=\left\lfloor\frac{4096}{4}\right\rfloor=1024
\ \text{escalares}.
\]

Um ciclo máximo de 64 estados comporta 65.536 dessas representações
escalares de comprimento fixo, antes do enquadramento. Em comparação com
UTF-8, o mapeamento escalar gasta mais bits com valores de um octeto e
menos com valores de quatro octetos. A comparação útil depende,
portanto, do repertório textual e do projeto do dicionário.

O RfD deixa o termo ``chave'' aberto a uma interpretação concreta:
dicionário público, dicionário privado provisionado, índice de estado
externo ou entrada de uma construção criptográfica. O mapeamento um a um
estabelece uma representação; seu eventual papel poderá ser escolhido na
discussão.

\hypertarget{estudo-reproduzuxedvel-e-exemplo-do-app}{%
\section{Estudo reproduzível e exemplo do
app}\label{estudo-reproduzuxedvel-e-exemplo-do-app}}

\hypertarget{o-estudo-de-16-estados}{%
\subsection{O estudo de 16 estados}\label{o-estudo-de-16-estados}}

A demonstração publicada contém deliberadamente 16 estados. Cada estado
usa o canvas atual de 256 × 256 pixels, 64 × 64 módulos efetivos e 64
entradas hexadecimais na paleta. A paleta ordenada é rotacionada entre
estados. O payload é sintético e reproduzível a partir da semente
20261104.

Seu orçamento bruto de dados é

\[
B_{16}=\frac{64\cdot64\cdot16\cdot6}{8}
=49.152\ \text{bytes}=48\ \text{KiB}.
\]

O decoder recupera exatamente o payload original. O relatório registra
os valores SHA-256 original e recuperado, o hash do GIF, versões do
ambiente, paletas hexadecimais, quantidade de estados e duração nominal.

\begin{longtable}[]{@{}lr@{}}
\toprule\noalign{}
Propriedade & Estudo medido \\
\midrule\noalign{}
\endhead
\bottomrule\noalign{}
\endlastfoot
Imagem armazenada & 256 × 256 pixels \\
Estado efetivo & 64 × 64 módulos \\
Frames & 16 \\
Entradas da paleta por estado & 64 \\
Ciclo nominal armazenado & 2.000 ms \\
Payload recuperado & 49.152 bytes \\
Tamanho do GIF codificado & 383.644 bytes \\
Recuperação & Exata, byte a byte \\
\end{longtable}

O tamanho do arquivo codificado é uma medida do contêiner. A contagem do
payload bruto vem da quantidade de dimensions independentes
selecionadas.

\hypertarget{implementauxe7uxe3o-e-verificauxe7uxf5es}{%
\subsection{Implementação e
verificações}\label{implementauxe7uxe3o-e-verificauxe7uxf5es}}

O código aceita octetos de payload em hexadecimal e entradas
hexadecimais de paleta. Escreve frames indexados completos, inclusive
estados repetidos, e lê os índices armazenados na ordem declarada pela
paleta. O decoder verifica que cada módulo efetivo contém 4 × 4 índices
idênticos.

A suíte de testes verifica:

\begin{itemize}
\tightlist
\item
  Recuperação exata para diferentes comprimentos de payload, incluindo
  octetos zero.
\item
  Preservação da ordem variável da paleta.
\item
  Registros de frame distintos para estados idênticos.
\item
  O limite do perfil de 64 frames e 192 KiB.
\item
  Cadência de um estado por minuto.
\item
  Limites de frames, paletas e entrada hexadecimal.
\item
  Rejeição de contêiner truncado e de módulo inconsistente.
\item
  Inversão dos endereços de quatro coordenadas e fronteiras dos
  escalares Unicode.
\end{itemize}

O visualizador local inicia em inglês, oferece português e pausa antes
de copiar o estado atual. Inclui cópia manual quando a operação direta
da área de transferência está indisponível. O código completo copiado no
exemplo contém 6.144 dígitos hexadecimais para um estado de 3 KiB.

O repositório fornece o gerador do estudo, resultados, visualizador,
fontes dos dois manuscritos e comandos de compilação.

\hypertarget{escopo-e-limitauxe7uxf5es}{%
\section{Escopo e limitações}\label{escopo-e-limitauxe7uxf5es}}

Esta seção concentra a interpretação e as fronteiras do RfD.

\textbf{Status de discussão.} Perfil, vocabulário e mapeamentos são
propostas de trabalho. Enquadramento, comprimento do código,
interoperabilidade, provisionamento e protocolo de verificação são temas
de discussão antes de uma especificação.

\textbf{Dimensões lógicas e contagem.} Os quatro eixos enumeram tuplas
candidatas. Uma dimension é selecionada por posição espacial e temporal.
As 16.777.216 entradas do catálogo e os 192 KiB de payload bruto
selecionado medem objetos distintos. ``cps'', quando empregado, designa
tuplas candidatas por segundo.

\textbf{Interface hexadecimal e contêiner.} Hexadecimal é a
representação pública dos valores da paleta e dos códigos. GIF89a
fornece sua serialização interna. Mudar a representação numérica não
cria entropia adicional. O adaptador atual usa frames completos e opacos
e paletas cujo tamanho é potência de dois, com até 64 entradas. O
comprimento original do payload é fornecido pelo chamador.

\textbf{Orçamentos e tempo.} Cabeçalhos, dicionários, dados de
autenticação, redundância e armazenamento do contêiner exigem orçamentos
separados. O exemplo de 8 fps descreve tempo nominal de reprodução; a
cadência dos estados do app é independente. O estudo não mede velocidade
de rede, vazão de modelos generativos nem captura física.

\textbf{Autenticação.} O exemplo do app contém estados sintéticos
públicos e demonstra seleção e cópia. Uma sequência finita repetida
exige um projeto escolhido de provisionamento e renovação para uso em
autenticação. Confidencialidade, vínculo com desafio, expiração,
tratamento de replay e limitação de tentativas pertencem a esse projeto.
\href{https://www.rfc-editor.org/rfc/rfc6238.html}{TOTP} e
\href{https://www.rfc-editor.org/rfc/rfc5116.html}{AEAD} são referências
para a discussão.

\textbf{Unicode e dicionários.} Valores escalares, caracteres
atribuídos, clusters de grafemas, octetos UTF-8, entradas ativas de
paleta e chaves possuem domínios diferentes. Os mapeamentos escalares
preservam a identidade escalar; normalização, ordem, multiplicidade e
armazenamento do dicionário exigem escolhas explícitas em um projeto de
texto ou chave.

\textbf{Evidência.} Os resultados registrados se aplicam ao caminho
atual de arquivo indexado exato. Sustentam as contagens e as alegações
de recuperação declaradas e não estabelecem a segurança de um protocolo
posterior.

\hypertarget{questuxf5es-para-discussuxe3o}{%
\section{Questões para discussão}\label{questuxf5es-para-discussuxe3o}}

\begin{enumerate}
\def\labelenumi{\arabic{enumi}.}
\tightlist
\item
  Quais aplicações se beneficiam do catálogo de quatro coordenadas:
  dados densos, eventos endereçados, seleção de estados ou consulta a
  dicionário?
\item
  Qual objeto hexadecimal deve ser chamado de chave e como deverá ser
  provisionado?
\item
  O texto deve usar octetos UTF-8, valores escalares em quatro
  dimensions ou dicionários adaptativos por estado?
\item
  Como associar o código selecionado a um serviço, desafio ou janela de
  aceitação?
\item
  Quais aplicações precisam de transporte rápido, cadência lenta no app
  ou avanço explícito pelo usuário?
\item
  O que deve renovar uma sequência após um ciclo ou após o consumo?
\item
  Quais campos são necessários antes da proposta de uma especificação?
\end{enumerate}

\hypertarget{autoria-e-assistuxeancia}{%
\section{Autoria e assistência}\label{autoria-e-assistuxeancia}}

\textbf{João Paulo De Luca Jr e Jorge Luiz Espanhol} são os coautores
deste RfD e da proposta de discussão associada. ChatGPT/Codex, OpenAI,
auxiliou na exposição matemática, tradução, implementação, execução e
preparação do documento. Os autores mantêm a responsabilidade pela
direção e pelas decisões de pesquisa.

\hypertarget{referuxeancias-tuxe9cnicas-com-links}{%
\section{Referências técnicas com
links}\label{referuxeancias-tuxe9cnicas-com-links}}

\begin{enumerate}
\def\labelenumi{\arabic{enumi}.}
\tightlist
\item
  \href{https://www.w3.org/Graphics/GIF/spec-gif89a.txt}{CompuServe,
  Graphics Interchange Format Version 89a}, reproduzido pelo W3C.
\item
  \href{https://pillow.readthedocs.io/en/stable/handbook/image-file-formats.html\#gif}{Pillow,
  suporte a arquivos GIF}.
\item
  \href{https://www.unicode.org/versions/Unicode18.0.0/core-spec/chapter-3/}{Unicode
  18.0, Capítulo 3: Conformidade, especialmente §3.9}.
\item
  \href{https://www.rfc-editor.org/rfc/rfc3629.html}{RFC 3629: UTF-8, a
  transformation format of ISO 10646}.
\item
  \href{https://www.rfc-editor.org/rfc/rfc6238.html}{RFC 6238: TOTP,
  Time-Based One-Time Password Algorithm}.
\item
  \href{https://www.rfc-editor.org/rfc/rfc5116.html}{RFC 5116: An
  Interface and Algorithms for Authenticated Encryption}.
\item
  \href{https://www.w3.org/TR/clipboard-apis/}{W3C, Clipboard API and
  events}.
\end{enumerate}

\end{document}
